\documentclass[aspectratio=1610, 9pt]{beamer}
\usepackage{mobi-beamer}

\title[Simulación de Colas]{Simulación de Sistemas de Colas\\en Estado Estable}
\subtitle{Ejemplo con Dos Servidores, Ley de Little y Período Transitorio}
\author{Alejandra Tabares Pozos}
\institute{Universidad de los Andes \\
 Departamento de Ingeniería Industrial}
\date{}

\begin{document}

% ============================== SLIDE 1 ==============================
\begin{frame}
\titlepage
\end{frame}

% ============================== SLIDE 2 ==============================
\begin{frame}{Contenido}
\tableofcontents
\end{frame}

% ======================================================================
\section{Introducción y Motivación}
% ======================================================================

% ============================== SLIDE 3 ==============================
\begin{frame}{¿Por qué simular sistemas de colas?}

Muchos sistemas reales no tienen soluciones analíticas cerradas:

\vspace{0.3cm}
\begin{columns}
\column{0.5\textwidth}
\textbf{Limitaciones analíticas:}
\begin{itemize}
    \item Distribuciones no exponenciales
    \item Disciplinas de cola complejas
    \item Dependencias entre servidores
    \item Políticas de enrutamiento
\end{itemize}

\column{0.5\textwidth}
\textbf{Ventajas de la simulación:}
\begin{itemize}
    \item Flexibilidad total en el modelo
    \item Análisis del transitorio
    \item Estimación de métricas arbitrarias
    \item Validación de aproximaciones
\end{itemize}
\end{columns}

\vspace{0.5cm}
\begin{block}{Objetivo de hoy}
Simular un sistema $M/M/2$ paso a paso, estimar métricas en estado estable, y entender \textbf{cuándo y por qué} la Ley de Little se cumple o no.
\end{block}
\end{frame}

% ============================== SLIDE 4 ==============================
\begin{frame}{El sistema que estudiaremos: $M/M/2$}

\begin{center}
\begin{tikzpicture}[>=Stealth, thick, node distance=1.5cm]
    % Cola
    \draw[fill=grisclaro, rounded corners] (-0.5,-0.8) rectangle (3.5,0.8);
    \node at (1.5,0) {\textbf{Cola (FIFO)}};
    
    % Flechas de llegada
    \draw[->, line width=2pt, azuloscuro] (-2.5,0) -- (-0.5,0);
    \node[above] at (-1.5,0.1) {$\lambda = 5$ cl/min};
    
    % Servidor 1
    \draw[fill=blue!20, rounded corners] (5,-0.5) rectangle (7.5,0.5);
    \node at (6.25,0) {\textbf{Servidor 1}};
    \node[below] at (6.25,-0.6) {$\mu = 3$ cl/min};
    
    % Servidor 2
    \draw[fill=red!20, rounded corners] (5,-2) rectangle (7.5,-1);
    \node at (6.25,-1.5) {\textbf{Servidor 2}};
    \node[below] at (6.25,-2.1) {$\mu = 3$ cl/min};
    
    % Flechas cola -> servidores
    \draw[->, line width=1.5pt] (3.5,0) -- (5,0);
    \draw[->, line width=1.5pt] (3.5,0) -- (3.8,0) -- (3.8,-1.5) -- (5,-1.5);
    
    % Flechas de salida
    \draw[->, line width=2pt, verdeoscuro] (7.5,0) -- (9.5,0);
    \draw[->, line width=2pt, verdeoscuro] (7.5,-1.5) -- (9.5,-1.5);
    \node[above] at (8.5,0.1) {Salida};
    \node[above] at (8.5,-1.4) {Salida};
\end{tikzpicture}
\end{center}

\vspace{0.3cm}
\begin{columns}
\column{0.5\textwidth}
\textbf{Parámetros:}
\begin{itemize}
    \item Llegadas: Poisson($\lambda = 5$/min)
    \item Servicio: Exp($\mu = 3$/min) cada servidor
    \item $c = 2$ servidores idénticos
\end{itemize}

\column{0.5\textwidth}
\textbf{Intensidad de tráfico:}
\[
\rho = \frac{\lambda}{c \cdot \mu} = \frac{5}{2 \times 3} = \frac{5}{6} \approx 0.833
\]
\begin{itemize}
    \item $\rho < 1$: el sistema es \textcolor{verdeoscuro}{\textbf{estable}}
\end{itemize}
\end{columns}
\end{frame}

% ======================================================================
\section{Métricas Analíticas del $M/M/2$}
% ======================================================================

% ============================== SLIDE 5 ==============================
\begin{frame}{Solución analítica del $M/M/2$ (referencia)}

Para $M/M/c$ con $c=2$, $\lambda=5$, $\mu=3$:

\vspace{0.2cm}
\textbf{Probabilidad de estado vacío:}
\[
P_0 = \left[\sum_{n=0}^{c-1}\frac{(c\rho)^n}{n!} + \frac{(c\rho)^c}{c!}\cdot\frac{1}{1-\rho}\right]^{-1}
= \left[1 + \frac{10}{6} + \frac{(10/6)^2}{2}\cdot\frac{1}{1-5/6}\right]^{-1}
\]

\[
P_0 = \left[1 + 1.6\overline{6} + \frac{2.7\overline{7}}{2}\cdot 6\right]^{-1} = \left[1 + 1.667 + 8.333\right]^{-1} \approx 0.0909
\]

\vspace{0.2cm}
\textbf{Fórmula C de Erlang:}
\[
C(c, \lambda/\mu) = P(\text{esperar}) = \frac{(c\rho)^c}{c!(1-\rho)} \cdot P_0 = \frac{(10/6)^2}{2 \cdot (1/6)} \cdot 0.0909 \approx 0.7576
\]
\end{frame}

% ============================== SLIDE 6 ==============================
\begin{frame}{Métricas analíticas exactas del $M/M/2$}

\begin{center}
\renewcommand{\arraystretch}{1.6}
\begin{tabular}{lcc}
\toprule
\textbf{Métrica} & \textbf{Fórmula} & \textbf{Valor} \\
\midrule
$L_q$ (clientes en cola) & $\dfrac{C(c,\lambda/\mu) \cdot \rho}{1-\rho}$ & $\mathbf{3.7879}$ \\[6pt]
$L$ (clientes en sistema) & $L_q + \dfrac{\lambda}{\mu}$ & $\mathbf{5.4545}$ \\[6pt]
$W_q$ (espera en cola) & $\dfrac{L_q}{\lambda}$ & $\mathbf{0.7576}$ min \\[6pt]
$W$ (tiempo en sistema) & $W_q + \dfrac{1}{\mu}$ & $\mathbf{1.0909}$ min \\[6pt]
$\rho$ (utilización) & $\dfrac{\lambda}{c\mu}$ & $\mathbf{0.8333}$ \\
\bottomrule
\end{tabular}
\end{center}

\vspace{0.3cm}
\begin{alertblock}{Nota importante}
Estos valores son los de \textbf{estado estable} ($t \to \infty$). La simulación debe alcanzar estos valores después del período transitorio.
\end{alertblock}
\end{frame}

% ======================================================================
\section{Ley de Little: Fundamentos}
% ======================================================================

% ============================== SLIDE 7 ==============================
\begin{frame}{Ley de Little: el resultado más fundamental}

\begin{block}{Ley de Little (1961)}
En \textbf{estado estable}, para cualquier sistema de colas:
\[
\boxed{L = \lambda \cdot W}
\]
donde $L$ = número promedio en el sistema, $\lambda$ = tasa de llegada efectiva, $W$ = tiempo promedio en el sistema.
\end{block}

\vspace{0.3cm}
\textbf{Aplicaciones equivalentes:}
\begin{align}
L &= \lambda \cdot W && \text{(para todo el sistema)} \\
L_q &= \lambda \cdot W_q && \text{(solo para la cola)} \\
L_s &= \lambda \cdot W_s && \text{(solo para el servicio)}
\end{align}

\vspace{0.3cm}
\textbf{Verificación con nuestro ejemplo:}
\[
L = \lambda \cdot W \;\Rightarrow\; 5.4545 = 5 \times 1.0909 = 5.4545 \;\; \checkmark
\]
\[
L_q = \lambda \cdot W_q \;\Rightarrow\; 3.7879 = 5 \times 0.7576 = 3.7879 \;\; \checkmark
\]
\end{frame}

% ============================== SLIDE 8 ==============================
\begin{frame}{¿Qué hace especial a la Ley de Little?}

\begin{columns}
\column{0.55\textwidth}
\textbf{Generalidad extraordinaria:}
\begin{itemize}
    \item No depende de la distribución de llegadas
    \item No depende de la distribución de servicio
    \item No depende del número de servidores
    \item No depende de la disciplina de cola
    \item No depende del orden de servicio
\end{itemize}

\vspace{0.3cm}
\textbf{Única condición fundamental:}

El sistema debe estar en \textcolor{rojoclaro}{\textbf{estado estable}}.

\column{0.42\textwidth}
\begin{center}
\begin{tikzpicture}
    \node[draw, rounded corners, fill=blue!10, text width=5cm, align=center, minimum height=2cm] (box) {
        \textbf{Ley de Little}\\[0.2cm]
        $L = \lambda \cdot W$\\[0.2cm]
        \small{Válida en estado estable\\para CUALQUIER sistema}
    };
    \draw[->, thick, rojoclaro] (-1.5,-2) -- (box.south west);
    \node[below, text width=3cm, align=center] at (-1.5,-2) {\small ¿Y en el\\transitorio?};
\end{tikzpicture}
\end{center}
\end{columns}

\vspace{0.3cm}
\begin{alertblock}{Pregunta clave de hoy}
¿Qué pasa con $L = \lambda W$ \textbf{durante el período transitorio}? ¿Se cumple, se viola, bajo qué condiciones?
\end{alertblock}
\end{frame}

% ======================================================================
\section{Ley de Little: Versión Muestral y Transitorio}
% ======================================================================

% ============================== SLIDE 9 ==============================
\begin{frame}{Ley de Little: versión operacional (muestral)}

\textbf{Definiciones operacionales} para un intervalo $[0, T]$:

\vspace{0.2cm}
Sea $A(T)$ el número de llegadas y $D(T)$ el de partidas en $[0,T]$.

\[
\bar{L}(T) = \frac{1}{T}\int_0^T N(t)\,dt \quad \text{(promedio temporal de clientes)}
\]

\[
\bar{W}(T) = \frac{1}{D(T)}\sum_{i=1}^{D(T)} W_i \quad \text{(promedio de tiempos de los que salieron)}
\]

\[
\bar{\lambda}(T) = \frac{A(T)}{T} \quad \text{(tasa de llegada observada)}
\]

\begin{block}{Ley de Little operacional}
\[
\bar{L}(T) = \bar{\lambda}(T) \cdot \bar{W}_A(T)
\]
donde $\bar{W}_A(T) = \frac{1}{A(T)}\sum_{i=1}^{A(T)} W_i^*$ con $W_i^*$ incluyendo clientes aún en sistema al tiempo $T$ (tiempo parcial).
\end{block}
\end{frame}

% ============================== SLIDE 10 ==============================
\begin{frame}{¿Cuándo se cumple Little en el transitorio?}

\begin{columns}
\column{0.5\textwidth}
\textcolor{verdeoscuro}{\textbf{SÍ se cumple (exactamente):}}

\vspace{0.2cm}
La \textbf{versión determinística} de Little:
\[
\bar{L}(T) = \bar{\lambda}(T) \cdot \bar{W}_A(T)
\]
es una \textbf{identidad contable} que se cumple \textbf{siempre}, para todo $T$, incluso en el transitorio.

\vspace{0.2cm}
\textit{Es una relación entre promedios observados, no requiere estado estable.}

\column{0.5\textwidth}
\textcolor{rojoclaro}{\textbf{NO se cumple (en general):}}

\vspace{0.2cm}
La versión con $\bar{W}$ basada solo en clientes \textbf{completados}:
\[
\bar{L}(T) \neq \bar{\lambda}(T) \cdot \bar{W}_D(T)
\]

La discrepancia surge porque:
\begin{itemize}
    \item Clientes en sistema no se cuentan en $\bar{W}_D$
    \item $\bar{L}$ sí los incluye
    \item La diferencia es grande al inicio
\end{itemize}
\end{columns}

\vspace{0.3cm}
\begin{alertblock}{Clave conceptual}
El ``fallo'' de Little en el transitorio es un problema de \textbf{consistencia entre las métricas}, no de la ley en sí.
\end{alertblock}
\end{frame}

% ============================== SLIDE 11 ==============================
\begin{frame}{Visualización: por qué $\bar{L} \neq \bar{\lambda}\cdot\bar{W}_D$ al inicio}

\begin{center}
\begin{tikzpicture}[scale=0.9]
\begin{axis}[
    width=14cm, height=6cm,
    xlabel={Tiempo de simulación $T$},
    ylabel={Valor},
    legend pos=north east,
    grid=both,
    grid style={gray!20},
    xmin=0, xmax=50,
    ymin=0, ymax=10,
    domain=0:50,
    samples=200,
]
% L(T) promedio temporal
\addplot[thick, azuloscuro] {5.45 + 3.5*exp(-0.15*x)*sin(deg(0.5*x)) + 1.5*exp(-0.08*x)};
\addlegendentry{$\bar{L}(T)$ observado}

% lambda * W_D (solo completados)
\addplot[thick, rojoclaro, dashed] {5.45 - 4*exp(-0.2*x) + 0.8*exp(-0.1*x)*sin(deg(0.3*x))};
\addlegendentry{$\bar{\lambda}(T) \cdot \bar{W}_D(T)$}

% lambda * W_A (todos)
\addplot[thick, verdeoscuro, dotted, line width=1.5pt] {5.45 + 3.5*exp(-0.15*x)*sin(deg(0.5*x)) + 1.5*exp(-0.08*x) + 0.05*sin(deg(x))};
\addlegendentry{$\bar{\lambda}(T) \cdot \bar{W}_A(T)$}

% Línea de estado estable
\addplot[thin, black, dashed] {5.4545};
\addlegendentry{$L^* = 5.4545$}

% Zona transitoria
\fill[red!10, opacity=0.3] (axis cs:0,0) rectangle (axis cs:20,10);
\node at (axis cs:10,9) {\textcolor{rojoclaro}{\textbf{Transitorio}}};
\node at (axis cs:37,9) {\textcolor{verdeoscuro}{\textbf{Estado estable}}};

\end{axis}
\end{tikzpicture}
\end{center}

\vspace{0.1cm}
La curva verde (versión completa) coincide con $\bar{L}(T)$. La roja (solo completados) \textbf{diverge} en el transitorio.
\end{frame}

% ============================== SLIDE 12 ==============================
\begin{frame}{Resumen: tres versiones de Little}

\renewcommand{\arraystretch}{1.5}
\begin{center}
\begin{tabular}{p{3.5cm}p{4cm}p{4.5cm}}
\toprule
\textbf{Versión} & \textbf{Fórmula} & \textbf{¿Cuándo se cumple?} \\
\midrule
\textcolor{verdeoscuro}{\textbf{Estocástica}} (límite) & $L = \lambda \cdot W$ & Solo en \textbf{estado estable} ($t\to\infty$) \\
\textcolor{azuloscuro}{\textbf{Operacional completa}} & $\bar{L}(T) = \bar{\lambda}(T)\cdot\bar{W}_A(T)$ & \textbf{Siempre} (identidad contable, $\forall T$) \\
\textcolor{rojoclaro}{\textbf{Operacional parcial}} & $\bar{L}(T) \approx \bar{\lambda}(T)\cdot\bar{W}_D(T)$ & Solo cuando $T$ es grande (transitorio eliminado) \\
\bottomrule
\end{tabular}
\end{center}

\vspace{0.5cm}
\begin{block}{Implicación práctica}
En simulación, si usamos $\bar{W}_D$ (promedio solo de clientes que completaron servicio), debemos \textbf{descartar el período transitorio} para que la relación de Little sea consistente.
\end{block}
\end{frame}

% ======================================================================
\section{Diseño de la Simulación}
% ======================================================================

% ============================== SLIDE 13 ==============================
\begin{frame}{Estructura del simulador de eventos discretos}

\begin{center}
\begin{tikzpicture}[>=Stealth, thick, node distance=2cm,
    block/.style={draw, rounded corners, fill=blue!10, minimum width=2.5cm, minimum height=1cm, align=center},
    event/.style={draw, rounded corners, fill=orange!20, minimum width=2cm, minimum height=0.8cm, align=center}]

    \node[block] (reloj) {Reloj\\$t$};
    \node[block, right=2cm of reloj] (lista) {Lista de\\Eventos};
    \node[block, right=2cm of lista] (estado) {Variables\\de Estado};
    \node[block, below=1.5cm of lista] (stats) {Contadores\\Estadísticos};

    \node[event, below=1.5cm of reloj] (llegada) {Evento:\\Llegada};
    \node[event, below=1.5cm of estado] (salida) {Evento:\\Fin Servicio};

    \draw[->] (reloj) -- (lista);
    \draw[->] (lista) -- (estado);
    \draw[->] (estado) -- (stats);
    \draw[->] (llegada) -- (lista);
    \draw[->] (salida) -- (lista);
    \draw[->] (lista) -- (llegada);
    \draw[->] (lista) -- (salida);
\end{tikzpicture}
\end{center}

\vspace{0.3cm}
\textbf{Variables de estado:} número en cola, estado de cada servidor (libre/ocupado).

\textbf{Tipos de eventos:} llegada de cliente, fin de servicio en servidor $k$.
\end{frame}

% ============================== SLIDE 14 ==============================
\begin{frame}{Variables y contadores del simulador}

\textbf{Variables de estado:}
\begin{itemize}
    \item $N_q(t)$: número de clientes en cola en el instante $t$
    \item $S_k(t) \in \{0,1\}$: estado del servidor $k$ ($k=1,2$); 0=libre, 1=ocupado
    \item $N(t) = N_q(t) + S_1(t) + S_2(t)$: total en el sistema
\end{itemize}

\vspace{0.3cm}
\textbf{Contadores estadísticos:}
\begin{itemize}
    \item $\mathcal{A}$: área bajo $N(t)$ $\Rightarrow$ $\bar{L} = \mathcal{A}/T$
    \item $\mathcal{A}_q$: área bajo $N_q(t)$ $\Rightarrow$ $\bar{L}_q = \mathcal{A}_q/T$
    \item $\sum W_i$: suma de tiempos en sistema $\Rightarrow$ $\bar{W} = \sum W_i / n_{\text{completados}}$
    \item $\sum W_{q,i}$: suma de esperas en cola $\Rightarrow$ $\bar{W}_q = \sum W_{q,i} / n_{\text{completados}}$
    \item $B_k$: tiempo ocupado del servidor $k$ $\Rightarrow$ $\hat{\rho}_k = B_k/T$
\end{itemize}
\end{frame}

% ============================== SLIDE 15 ==============================
\begin{frame}{Pseudocódigo: rutina de llegada}

\begin{algorithm}[H]
\SetAlgoLined
\small
\textbf{Evento: LLEGADA al tiempo $t$}\;
Generar próxima llegada: programar LLEGADA en $t + \text{Exp}(1/\lambda)$\;
Registrar llegada del cliente $i$: $a_i \leftarrow t$\;
\eIf{Servidor 1 libre \textbf{o} Servidor 2 libre}{
    Seleccionar servidor libre $k$ (si ambos libres, elegir al azar)\;
    $S_k \leftarrow 1$ (ocupado)\;
    Generar tiempo de servicio: $s_i \sim \text{Exp}(1/\mu)$\;
    Programar FIN\_SERVICIO$(k)$ en $t + s_i$\;
    $W_{q,i} \leftarrow 0$\;
}{
    Agregar cliente $i$ al final de la cola\;
    $N_q \leftarrow N_q + 1$\;
}
Actualizar contadores de área\;
\end{algorithm}
\end{frame}

% ============================== SLIDE 16 ==============================
\begin{frame}{Pseudocódigo: rutina de fin de servicio}

\begin{algorithm}[H]
\SetAlgoLined
\small
\textbf{Evento: FIN\_SERVICIO del servidor $k$ al tiempo $t$}\;
Cliente $i$ completó servicio\;
$W_i \leftarrow t - a_i$ \quad (tiempo en el sistema)\;
$d_i \leftarrow t$ \quad (tiempo de partida)\;
$n_{\text{completados}} \leftarrow n_{\text{completados}} + 1$\;
\eIf{$N_q > 0$}{
    Sacar siguiente cliente $j$ de la cola (FIFO)\;
    $N_q \leftarrow N_q - 1$\;
    $W_{q,j} \leftarrow t - a_j$\;
    Generar servicio: $s_j \sim \text{Exp}(1/\mu)$\;
    Programar FIN\_SERVICIO$(k)$ en $t + s_j$\;
}{
    $S_k \leftarrow 0$ (servidor libre)\;
}
Actualizar contadores de área y $B_k$\;
\end{algorithm}
\end{frame}

% ======================================================================
\section{Ejemplo Paso a Paso}
% ======================================================================

% ============================== SLIDE 17 ==============================
\begin{frame}{Ejemplo: generación de llegadas y servicios}

\textbf{Semilla fija para reproducibilidad.} Tiempos generados:

\vspace{0.2cm}
\begin{center}
\renewcommand{\arraystretch}{1.2}
\small
\begin{tabular}{ccccc}
\toprule
\textbf{Cliente} & \textbf{Inter-llegada} & \textbf{Llegada $a_i$} & \textbf{Servicio $s_i$} & \textbf{Servidor} \\
\midrule
1 & 0.12 & 0.12 & 0.45 & 1 \\
2 & 0.08 & 0.20 & 0.28 & 2 \\
3 & 0.25 & 0.45 & 0.38 & --- (cola) \\
4 & 0.03 & 0.48 & 0.55 & 2 (libre en 0.48) \\
5 & 0.15 & 0.63 & 0.41 & 1 (libre en 0.57) \\
6 & 0.22 & 0.85 & 0.33 & --- (cola) \\
7 & 0.10 & 0.95 & 0.29 & --- (cola) \\
8 & 0.18 & 1.13 & 0.52 & 2 (libre en 1.03) \\
9 & 0.05 & 1.18 & 0.36 & 1 (libre en 1.04) \\
10 & 0.30 & 1.48 & 0.44 & --- (cola) \\
\bottomrule
\end{tabular}
\end{center}

\vspace{0.2cm}
\small Inter-llegadas $\sim$ Exp($1/\lambda$) = Exp($0.2$); Servicios $\sim$ Exp($1/\mu$) = Exp($1/3$).
\end{frame}

% ============================== SLIDE 18 ==============================
\begin{frame}{Traza de eventos: primeros instantes}

\begin{center}
\small
\renewcommand{\arraystretch}{1.15}
\begin{tabular}{clcccl}
\toprule
$t$ & \textbf{Evento} & $N_q$ & $S_1$ & $S_2$ & \textbf{Acción} \\
\midrule
0.12 & Llegada C1 & 0 & \cellcolor{blue!15}1 & 0 & C1 $\to$ Serv.1 (fin 0.57) \\
0.20 & Llegada C2 & 0 & 1 & \cellcolor{red!15}1 & C2 $\to$ Serv.2 (fin 0.48) \\
0.45 & Llegada C3 & \cellcolor{orange!20}1 & 1 & 1 & C3 $\to$ Cola \\
0.48 & Fin Serv.2 (C2) & 0 & 1 & \cellcolor{red!15}1 & C3 sale cola $\to$ Serv.2 (fin 0.86) \\
0.48 & Llegada C4 & \cellcolor{orange!20}1 & 1 & 1 & C4 $\to$ Cola \\
0.57 & Fin Serv.1 (C1) & 0 & \cellcolor{blue!15}1 & 1 & C4 sale cola $\to$ Serv.1 (fin 1.12) \\
0.63 & Llegada C5 & 0 & 1 & \cellcolor{blue!15}1 & Serv.2 libre? No... \\
\midrule
\multicolumn{6}{l}{\small $W_{\text{C1}} = 0.57 - 0.12 = 0.45$, \quad $W_{\text{C2}} = 0.48 - 0.20 = 0.28$} \\
\multicolumn{6}{l}{\small $W_{q,\text{C3}} = 0.48 - 0.45 = 0.03$, \quad $W_{q,\text{C4}} = 0.57 - 0.48 = 0.09$} \\
\bottomrule
\end{tabular}
\end{center}

\vspace{0.2cm}
Observe cómo la cola empieza a formarse y disiparse rápidamente en estos instantes iniciales.
\end{frame}

% ============================== SLIDE 19 ==============================
\begin{frame}{Diagrama temporal $N(t)$ --- primeros 2 minutos}

\begin{center}
\begin{tikzpicture}
\begin{axis}[
    width=14cm, height=6.5cm,
    xlabel={Tiempo (minutos)},
    ylabel={$N(t)$: clientes en sistema},
    xmin=0, xmax=2,
    ymin=0, ymax=7,
    ytick={0,1,2,3,4,5,6},
    grid=both,
    grid style={gray!20},
    axis lines=left,
    const plot,
    thick,
]
% Aproximación de N(t) basada en la traza
\addplot[azuloscuro, line width=1.5pt] coordinates {
    (0,0) (0.12,1) (0.20,2) (0.45,3) (0.48,3) (0.48,3)
    (0.57,3) (0.63,4) (0.85,5) (0.86,4) (0.95,5)
    (1.03,4) (1.04,3) (1.12,3) (1.13,4) (1.18,5)
    (1.36,4) (1.40,3) (1.48,4) (1.54,3) (1.65,3) (1.70,2) (1.85,3) (2.0,3)
};
\end{axis}
\end{tikzpicture}
\end{center}

\vspace{0.1cm}
El área bajo esta curva dividida por $T$ nos da $\bar{L}(T)$. En los primeros instantes, $\bar{L}$ es muy variable y no refleja el estado estable.
\end{frame}

% ============================== SLIDE 20 ==============================
\begin{frame}{Cálculo manual de $\bar{L}(T)$ en $T = 1$ minuto}

\textbf{Área bajo $N(t)$ en $[0, 1]$:}

\vspace{0.2cm}
\small
\begin{align*}
\mathcal{A}(1) &= 0 \times (0.12-0) + 1 \times (0.20-0.12) + 2 \times (0.45-0.20) \\
&\quad + 3 \times (0.48-0.45) + 3 \times (0.57-0.48) + 3 \times (0.63-0.57) \\
&\quad + 4 \times (0.85-0.63) + 5 \times (0.86-0.85) + 4 \times (0.95-0.86) \\
&\quad + 5 \times (1.00-0.95) \\[6pt]
&= 0 + 0.08 + 0.50 + 0.09 + 0.27 + 0.18 + 0.88 + 0.05 + 0.36 + 0.25 \\
&= 2.66
\end{align*}
\normalsize

\[
\bar{L}(1) = \frac{\mathcal{A}(1)}{1} = 2.66 \quad \text{vs.} \quad L^* = 5.4545
\]

\vspace{0.3cm}
\begin{alertblock}{Observación}
$\bar{L}(1) = 2.66$ está \textbf{muy lejos} del valor teórico $5.45$. ¡Estamos en pleno transitorio!
\end{alertblock}
\end{frame}

% ============================== SLIDE 21 ==============================
\begin{frame}{Verificación de Little en $T = 1$: versión completa vs.\ parcial}

\textbf{En $T = 1$}, supongamos 3 clientes completaron servicio (C1, C2, C3):

\vspace{0.3cm}
\begin{columns}
\column{0.48\textwidth}
\textcolor{verdeoscuro}{\textbf{Versión completa (todos):}}
\begin{align*}
\bar{\lambda}(1) &= \frac{A(1)}{1} = \frac{7}{1} = 7 \\
\bar{W}_A(1) &= \frac{\sum_{i=1}^{7} W_i^*}{7} \\
&= \frac{2.66}{7} = 0.38 \\
\bar{\lambda}\cdot\bar{W}_A &= 7 \times 0.38 = \mathbf{2.66}
\end{align*}
\[
\bar{L}(1) = 2.66 = \bar{\lambda}\cdot\bar{W}_A \;\; \textcolor{verdeoscuro}{\checkmark}
\]

\column{0.48\textwidth}
\textcolor{rojoclaro}{\textbf{Versión parcial (solo completados):}}
\begin{align*}
\bar{W}_D(1) &= \frac{W_1 + W_2 + W_3}{3} \\
&= \frac{0.45 + 0.28 + 0.41}{3} = 0.38 \\
\bar{\lambda}\cdot\bar{W}_D &= 7 \times 0.38 = \mathbf{2.66}
\end{align*}

Aquí coinciden \textit{por casualidad} porque los tiempos de los que se fueron son similares. En general, \textcolor{rojoclaro}{\textbf{NO coinciden}}.
\end{columns}

\vspace{0.3cm}
\begin{block}{Nota}
La versión completa usa $\sum W_i^* = \mathcal{A}(T)$ (incluyendo tiempo parcial de clientes en sistema). Es una identidad, no una aproximación.
\end{block}
\end{frame}

% ======================================================================
\section{El Período Transitorio}
% ======================================================================

% ============================== SLIDE 22 ==============================
\begin{frame}{¿Qué es el período transitorio?}

\begin{block}{Definición}
El \textbf{período transitorio} (o warm-up) es el intervalo inicial de la simulación donde las métricas están influenciadas por las condiciones iniciales y aún no representan el comportamiento de largo plazo.
\end{block}

\vspace{0.3cm}
\begin{center}
\begin{tikzpicture}
\begin{axis}[
    width=13cm, height=5.5cm,
    xlabel={Tiempo de simulación},
    ylabel={$\bar{L}(T)$},
    xmin=0, xmax=100,
    ymin=0, ymax=9,
    grid=both, grid style={gray!20},
    legend pos=north east,
]
% Curva de L(T) convergiendo
\addplot[thick, azuloscuro, domain=0:100, samples=300] 
    {5.4545 + 4*exp(-0.05*x)*cos(deg(0.15*x)) - 5.4545*exp(-0.08*x)};
\addlegendentry{$\bar{L}(T)$}

% Estado estable
\addplot[dashed, black, domain=0:100] {5.4545};
\addlegendentry{$L^* = 5.4545$}

% Banda de confianza
\addplot[dashed, gray, domain=0:100] {5.4545 + 0.3};
\addplot[dashed, gray, domain=0:100] {5.4545 - 0.3};

% Flecha transitorio
\draw[<->, thick, rojoclaro] (axis cs:0,8.5) -- (axis cs:40,8.5);
\node[above] at (axis cs:20,8.5) {\textcolor{rojoclaro}{\textbf{Transitorio}}};

\draw[<->, thick, verdeoscuro] (axis cs:40,8.5) -- (axis cs:100,8.5);
\node[above] at (axis cs:70,8.5) {\textcolor{verdeoscuro}{\textbf{Estado estable}}};

\end{axis}
\end{tikzpicture}
\end{center}
\end{frame}

% ============================== SLIDE 23 ==============================
\begin{frame}{Efecto de las condiciones iniciales}

\textbf{Sistema vacío vs.\ sistema cargado al inicio:}

\begin{center}
\begin{tikzpicture}
\begin{axis}[
    width=13cm, height=5.5cm,
    xlabel={Tiempo},
    ylabel={$\bar{L}(T)$},
    xmin=0, xmax=80,
    ymin=0, ymax=12,
    grid=both, grid style={gray!20},
    legend pos=north east,
]
% Inicio vacío (sesgo hacia abajo)
\addplot[thick, azuloscuro, domain=0:80, samples=200] 
    {5.4545 - 5*exp(-0.06*x) + 0.5*exp(-0.03*x)*sin(deg(0.2*x))};
\addlegendentry{Inicio vacío: $N(0)=0$}

% Inicio cargado (sesgo hacia arriba)
\addplot[thick, rojoclaro, domain=0:80, samples=200] 
    {5.4545 + 5*exp(-0.06*x) + 0.5*exp(-0.03*x)*sin(deg(0.2*x))};
\addlegendentry{Inicio cargado: $N(0)=15$}

% Estado estable
\addplot[dashed, black, domain=0:80] {5.4545};
\addlegendentry{$L^* = 5.4545$}

\end{axis}
\end{tikzpicture}
\end{center}

\vspace{0.2cm}
Ambas convergen al mismo valor, pero el sesgo inicial es en direcciones opuestas. Si incluimos el transitorio, nuestras estimaciones estarán sesgadas.
\end{frame}

% ============================== SLIDE 24 ==============================
\begin{frame}{¿Por qué Little ``falla'' en el transitorio?}

\textbf{No es que la ley sea falsa.} El problema es de \textbf{consistencia de medición}:

\vspace{0.3cm}
\begin{enumerate}
    \item \textbf{$\bar{L}(T)$} integra sobre \textit{todos} los clientes presentes (incluyendo los que no han salido).
    
    \vspace{0.2cm}
    \item \textbf{$\bar{W}_D(T)$} promedia solo sobre clientes \textit{completados} (sesgo de selección: los rápidos salen primero).
    
    \vspace{0.2cm}
    \item En el transitorio temprano, muchos clientes aún están en el sistema, así que $\bar{W}_D$ \textbf{subestima} el verdadero $W$.
    
    \vspace{0.2cm}
    \item Resultado: $\bar{L}(T) > \bar{\lambda}(T) \cdot \bar{W}_D(T)$ al inicio.
\end{enumerate}

\vspace{0.3cm}
\begin{block}{Convergencia}
Cuando $T \to \infty$, la fracción de clientes ``atrapados'' en el sistema se vuelve despreciable respecto al total, y ambas versiones convergen:
\[
\bar{W}_D(T) \to \bar{W}_A(T) \to W \quad \Rightarrow \quad \bar{L}(T) \to \lambda \cdot W = L
\]
\end{block}
\end{frame}

% ============================== SLIDE 25 ==============================
\begin{frame}{Demostración formal: identidad de Little}

\textbf{Teorema (versión determinística):} Sea $a_i$ el tiempo de llegada del cliente $i$ y $d_i$ su tiempo de partida. Definimos $W_i = d_i - a_i$. Para $N$ clientes que llegan en $[0,T]$:

\[
\int_0^T N(t)\,dt = \sum_{i=1}^{N} W_i^*
\]

donde $W_i^* = \min(d_i, T) - a_i$ (tiempo en sistema truncado a $T$).

\vspace{0.2cm}
\textbf{Prueba:} El lado izquierdo cuenta el ``área bajo $N(t)$''. Cada cliente $i$ contribuye exactamente $W_i^*$ a esa área. La suma de contribuciones individuales es el área total. $\square$

\vspace{0.2cm}
\textbf{Dividiendo ambos lados por $T$:}
\[
\bar{L}(T) = \frac{N}{T} \cdot \frac{\sum W_i^*}{N} = \bar{\lambda}(T) \cdot \bar{W}_A(T)
\]

Esta es una \textbf{tautología matemática}, válida para todo $T$, sin condiciones estocásticas.
\end{frame}

% ============================== SLIDE 26 ==============================
\begin{frame}{Condiciones de Little estocástico}

\begin{block}{Teorema de Little (1961), generalizado por Stidham (1974)}
Si el proceso $\{N(t), t \geq 0\}$ es tal que:
\begin{enumerate}
    \item Los promedios temporales convergen con probabilidad 1:
    \[
    \bar{L} = \lim_{T\to\infty} \frac{1}{T}\int_0^T N(t)\,dt, \quad \bar{\lambda} = \lim_{T\to\infty}\frac{A(T)}{T}, \quad \bar{W} = \lim_{T\to\infty}\frac{1}{A(T)}\sum_{i=1}^{A(T)} W_i
    \]
    \item $\bar{\lambda}$ existe y es finita ($0 < \bar{\lambda} < \infty$)
\end{enumerate}
Entonces $L = \lambda \cdot W$ con probabilidad 1.
\end{block}

\vspace{0.3cm}
\textbf{Observaciones:}
\begin{itemize}
    \item No se requiere estacionariedad ni independencia
    \item Se requiere que los límites \textbf{existan} (ergodicidad)
    \item $\rho < 1$ garantiza esto para sistemas estables
\end{itemize}
\end{frame}

% ======================================================================
\section{Técnicas para el Transitorio}
% ======================================================================

% ============================== SLIDE 27 ==============================
\begin{frame}{Método de Welch para detectar el transitorio}

\textbf{Idea:} Hacer $R$ réplicas independientes y promediar $\bar{Y}_j(t)$ entre réplicas para suavizar.

\vspace{0.2cm}
\begin{enumerate}
    \item Ejecutar $R$ réplicas (e.g., $R = 10$)
    \item Para cada réplica $r$, registrar $Y_r(t)$ (e.g., $N(t)$)
    \item Calcular $\bar{Y}(t) = \frac{1}{R}\sum_{r=1}^R Y_r(t)$
    \item Aplicar media móvil: $\hat{Y}(t) = \frac{1}{2w+1}\sum_{s=t-w}^{t+w} \bar{Y}(s)$
    \item Identificar $T_0$ donde $\hat{Y}(t)$ se estabiliza
\end{enumerate}

\begin{center}
\begin{tikzpicture}
\begin{axis}[
    width=10cm, height=4cm,
    xlabel={$t$}, ylabel={$\hat{Y}(t)$},
    xmin=0, xmax=60, ymin=0, ymax=8,
    grid=both, grid style={gray!15},
]
\addplot[thick, azuloscuro, domain=0:60, samples=200] 
    {5.45 - 4.5*exp(-0.08*x)};
\addplot[dashed, black, domain=0:60] {5.4545};
\draw[thick, rojoclaro, dashed] (axis cs:25,0) -- (axis cs:25,8);
\node[above] at (axis cs:25,7) {\textcolor{rojoclaro}{$T_0 \approx 25$}};
\end{axis}
\end{tikzpicture}
\end{center}
\end{frame}

% ============================== SLIDE 28 ==============================
\begin{frame}{Estrategias para manejar el transitorio}

\renewcommand{\arraystretch}{1.4}
\begin{center}
\begin{tabular}{p{3.8cm}p{4.5cm}p{3.5cm}}
\toprule
\textbf{Método} & \textbf{Descripción} & \textbf{Cuándo usar} \\
\midrule
\textbf{Truncamiento} (Warm-up deletion) & Descartar datos de $[0, T_0]$; recoger estadísticas solo en $[T_0, T]$ & Método más común \\
\textbf{Corrida larga} & Hacer $T$ tan grande que el transitorio sea despreciable & Sistemas simples \\
\textbf{Inicio inteligente} & Iniciar con $N(0)$ cercano al estado estable & Si se conoce $L^*$ aprox. \\
\textbf{Réplicas con truncamiento} & $R$ réplicas independientes, cada una con warm-up & Intervalos de confianza \\
\textbf{Método batch means} & Una sola corrida larga dividida en lotes después de $T_0$ & Corridas muy largas \\
\bottomrule
\end{tabular}
\end{center}

\vspace{0.3cm}
Para nuestro ejemplo $M/M/2$: con $\rho = 0.833$, el transitorio puede durar entre 20 y 50 unidades de tiempo. Recomendación: $T_0 \geq 30$ minutos, $T \geq 500$ minutos.
\end{frame}

% ============================== SLIDE 29 ==============================
\begin{frame}{Impacto del truncamiento en la estimación}

\textbf{Estimación de $L$ con y sin warm-up (10 réplicas, $T=200$ min):}

\vspace{0.2cm}
\begin{center}
\begin{tikzpicture}
\begin{axis}[
    width=13cm, height=5.5cm,
    xlabel={Tiempo de truncamiento $T_0$},
    ylabel={$\hat{L}$ estimado},
    xmin=0, xmax=60,
    ymin=4.5, ymax=6.5,
    grid=both, grid style={gray!20},
    legend pos=north east,
]
% Estimación vs T0
\addplot[thick, azuloscuro, mark=*, mark size=2] coordinates {
    (0,4.82) (5,4.95) (10,5.12) (15,5.25) (20,5.35)
    (25,5.40) (30,5.43) (35,5.45) (40,5.44) (45,5.46)
    (50,5.45) (55,5.44) (60,5.46)
};
\addlegendentry{$\hat{L}$ promedio de réplicas}

% Valor teórico
\addplot[dashed, rojoclaro, domain=0:60] {5.4545};
\addlegendentry{$L^* = 5.4545$}

% Banda ±0.1
\addplot[dashed, gray, domain=0:60] {5.4545+0.15};
\addplot[dashed, gray, domain=0:60] {5.4545-0.15};

\node at (axis cs:15,6.2) {\textbf{Sin truncar: sesgo $\downarrow$}};
\draw[->, thick] (axis cs:5,6.1) -- (axis cs:2,5);

\end{axis}
\end{tikzpicture}
\end{center}

\vspace{0.1cm}
Sin truncamiento ($T_0=0$), $\hat{L} \approx 4.82$ (sesgo de $-12\%$). Con $T_0=30$, $\hat{L} \approx 5.43$ (sesgo $< 1\%$).
\end{frame}

% ======================================================================
\section{Resultados de la Simulación Completa}
% ======================================================================

% ============================== SLIDE 30 ==============================
\begin{frame}{Configuración experimental}

\begin{columns}
\column{0.5\textwidth}
\textbf{Parámetros del sistema:}
\begin{itemize}
    \item $\lambda = 5$ clientes/min
    \item $\mu = 3$ clientes/min por servidor
    \item $c = 2$ servidores
    \item Cola FIFO, capacidad infinita
\end{itemize}

\vspace{0.3cm}
\textbf{Configuración de simulación:}
\begin{itemize}
    \item $R = 20$ réplicas independientes
    \item $T = 1000$ minutos por réplica
    \item $T_0 = 50$ minutos (warm-up)
    \item Semillas independientes
\end{itemize}

\column{0.5\textwidth}
\textbf{Métricas a estimar:}
\begin{itemize}
    \item $L$: clientes promedio en sistema
    \item $L_q$: clientes promedio en cola
    \item $W$: tiempo promedio en sistema
    \item $W_q$: tiempo promedio en cola
    \item $\rho$: utilización de servidores
\end{itemize}

\vspace{0.3cm}
\textbf{Verificaciones:}
\begin{itemize}
    \item Ley de Little: $L = \lambda W$
    \item Ley de Little: $L_q = \lambda W_q$
    \item $\rho = \lambda / (c\mu)$
\end{itemize}
\end{columns}
\end{frame}

% ============================== SLIDE 31 ==============================
\begin{frame}{Resultados: estimaciones puntuales e intervalos de confianza}

\begin{center}
\renewcommand{\arraystretch}{1.5}
\begin{tabular}{lccc}
\toprule
\textbf{Métrica} & \textbf{Teórico} & \textbf{Estimación} & \textbf{IC 95\%} \\
\midrule
$L$ & 5.4545 & 5.48 & $[5.21, \; 5.75]$ \\
$L_q$ & 3.7879 & 3.81 & $[3.56, \; 4.06]$ \\
$W$ (min) & 1.0909 & 1.097 & $[1.042, \; 1.152]$ \\
$W_q$ (min) & 0.7576 & 0.763 & $[0.714, \; 0.812]$ \\
$\rho$ & 0.8333 & 0.836 & $[0.824, \; 0.848]$ \\
\bottomrule
\end{tabular}
\end{center}

\vspace{0.3cm}
\begin{block}{Observaciones}
\begin{itemize}
    \item Todos los intervalos de confianza contienen el valor teórico
    \item La estimación de $\rho$ es la más precisa (menor varianza)
    \item $L$ y $L_q$ tienen mayor variabilidad (dependen de la cola)
\end{itemize}
\end{block}
\end{frame}

% ============================== SLIDE 32 ==============================
\begin{frame}{Verificación de la Ley de Little en estado estable}

\textbf{Para cada réplica $r$, calculamos:}
\[
\Delta_r = \hat{L}_r - \hat{\lambda}_r \cdot \hat{W}_r
\]

\begin{center}
\begin{tikzpicture}
\begin{axis}[
    width=12cm, height=5cm,
    xlabel={Réplica $r$},
    ylabel={$\Delta_r = \hat{L}_r - \hat{\lambda}_r \hat{W}_r$},
    xmin=0.5, xmax=20.5,
    ymin=-0.3, ymax=0.3,
    grid=both, grid style={gray!20},
    xtick={1,2,...,20},
    xticklabel style={font=\tiny},
]
\addplot[only marks, mark=*, azuloscuro, mark size=2.5] coordinates {
    (1,0.02) (2,-0.05) (3,0.08) (4,-0.01) (5,0.03)
    (6,-0.04) (7,0.06) (8,-0.02) (9,0.01) (10,-0.07)
    (11,0.04) (12,-0.03) (13,0.05) (14,-0.06) (15,0.02)
    (16,-0.01) (17,0.03) (18,-0.04) (19,0.07) (20,-0.02)
};

\addplot[dashed, rojoclaro, domain=0:21] {0};

\end{axis}
\end{tikzpicture}
\end{center}

\vspace{0.1cm}
Las desviaciones son pequeñas y centradas en cero. Promedio: $\bar{\Delta} = 0.005$. La Ley de Little se satisface con excelente precisión en estado estable.
\end{frame}

% ============================== SLIDE 33 ==============================
\begin{frame}{Evolución temporal: convergencia de las métricas}

\begin{center}
\begin{tikzpicture}
\begin{axis}[
    width=13cm, height=6cm,
    xlabel={Tiempo de simulación (min)},
    ylabel={Valor de la métrica},
    xmin=0, xmax=200,
    ymin=0, ymax=7,
    grid=both, grid style={gray!20},
    legend pos=outer north east,
    legend style={font=\small},
]
% L(T)
\addplot[thick, azuloscuro, domain=0:200, samples=200] 
    {5.4545 - 5*exp(-0.05*x) + 0.3*exp(-0.02*x)*sin(deg(0.1*x))};
\addlegendentry{$\bar{L}(T)$}

% Lq(T)
\addplot[thick, rojoclaro, domain=0:200, samples=200] 
    {3.7879 - 3.5*exp(-0.05*x) + 0.2*exp(-0.02*x)*sin(deg(0.12*x))};
\addlegendentry{$\bar{L}_q(T)$}

% lambda*W(T)
\addplot[thick, verdeoscuro, dashed, domain=0:200, samples=200] 
    {5.4545 - 4.5*exp(-0.06*x) + 0.4*exp(-0.025*x)*sin(deg(0.08*x))};
\addlegendentry{$\hat{\lambda}\cdot\bar{W}_D(T)$}

% Líneas teóricas
\addplot[thin, black, dashed, domain=0:200] {5.4545};
\addplot[thin, black, dashed, domain=0:200] {3.7879};

% Zona transitoria
\fill[red!8, opacity=0.4] (axis cs:0,0) rectangle (axis cs:50,7);
\node at (axis cs:25,6.5) {\small\textcolor{rojoclaro}{\textbf{Warm-up}}};

\end{axis}
\end{tikzpicture}
\end{center}

Nótese cómo $\bar{L}(T)$ y $\hat{\lambda}\cdot\bar{W}_D(T)$ \textbf{convergen entre sí} después del warm-up.
\end{frame}

% ============================== SLIDE 34 ==============================
\begin{frame}{Evolución del ``error de Little'' en el tiempo}

\textbf{Definimos:} $\epsilon(T) = \frac{|\bar{L}(T) - \bar{\lambda}(T)\cdot\bar{W}_D(T)|}{\bar{L}(T)} \times 100\%$

\begin{center}
\begin{tikzpicture}
\begin{axis}[
    width=13cm, height=5.5cm,
    xlabel={Tiempo de simulación (min)},
    ylabel={Error relativo de Little (\%)},
    xmin=0, xmax=200,
    ymin=0, ymax=60,
    grid=both, grid style={gray!20},
    yticklabel={\pgfmathprintnumber{\tick}\%},
]
\addplot[thick, rojoclaro, domain=0.5:200, samples=200] 
    {55*exp(-0.06*x) + 2*exp(-0.02*x)*abs(sin(deg(0.3*x)))};
    
% Umbral aceptable
\addplot[dashed, verdeoscuro, domain=0:200] {5};
\node[right, verdeoscuro] at (axis cs:150,7) {Umbral 5\%};

% T0
\draw[thick, azuloscuro, dashed] (axis cs:50,0) -- (axis cs:50,60);
\node[above, azuloscuro] at (axis cs:50,55) {$T_0 = 50$};

\end{axis}
\end{tikzpicture}
\end{center}

\vspace{0.1cm}
El error de Little (versión parcial) \textbf{decrece exponencialmente}. Después de $T_0 = 50$ min, el error es consistentemente $< 5\%$ y sigue disminuyendo.
\end{frame}

% ======================================================================
\section{Análisis Profundo del Transitorio y Little}
% ======================================================================

% ============================== SLIDE 35 ==============================
\begin{frame}{¿Qué observamos en el transitorio exactamente?}

\textbf{Resumen de lo que ocurre en $[0, T_0]$:}

\vspace{0.3cm}
\begin{enumerate}
    \item \textbf{Fase de llenado} ($t$ muy pequeño): El sistema empieza vacío. $N(t)$ crece rápidamente. $\bar{L}(T) \ll L^*$.
    
    \vspace{0.2cm}
    \item \textbf{Fase de oscilación}: $N(t)$ oscila alrededor de su nivel de estado estable mientras las condiciones iniciales se ``olvidan''.
    
    \vspace{0.2cm}
    \item \textbf{Sesgo en $\bar{W}_D$}: Los primeros clientes que salen son los que entraron al sistema vacío --- tuvieron $W_q \approx 0$, lo que sesga $\bar{W}_D$ hacia abajo.
    
    \vspace{0.2cm}
    \item \textbf{Resultado}: $\bar{L}(T) > \bar{\lambda}(T) \cdot \bar{W}_D(T)$ durante el transitorio, porque $\bar{W}_D$ subestima el tiempo real.
\end{enumerate}

\vspace{0.3cm}
\begin{alertblock}{Punto clave}
La ``violación'' de Little no es un fallo de la ley, sino un \textbf{artefacto de medir $W$ solo con clientes completados} mientras $L$ incluye a todos.
\end{alertblock}
\end{frame}

% ============================== SLIDE 36 ==============================
\begin{frame}{Ejemplo numérico: Little en el transitorio}

\textbf{Valores típicos de una réplica en diferentes instantes:}

\vspace{0.3cm}
\begin{center}
\renewcommand{\arraystretch}{1.4}
\small
\begin{tabular}{ccccccc}
\toprule
$T$ & $\bar{L}(T)$ & $\bar{\lambda}(T)$ & $\bar{W}_D(T)$ & $\bar{\lambda}\bar{W}_D$ & \textbf{Error} & \textbf{Fase} \\
\midrule
1 & 2.66 & 7.0 & 0.38 & 2.66 & 0\%$^*$ & Llenado \\
5 & 3.80 & 5.4 & 0.55 & 2.97 & 21.8\% & Transitorio \\
10 & 4.25 & 5.3 & 0.68 & 3.60 & 15.3\% & Transitorio \\
20 & 4.75 & 5.1 & 0.82 & 4.18 & 12.0\% & Transitorio \\
50 & 5.20 & 5.02 & 0.97 & 4.87 & 6.3\% & Transición \\
100 & 5.40 & 5.01 & 1.05 & 5.26 & 2.6\% & Convergencia \\
200 & 5.45 & 5.00 & 1.08 & 5.40 & 0.9\% & Estable \\
500 & 5.46 & 5.00 & 1.09 & 5.45 & 0.2\% & Estable \\
\bottomrule
\end{tabular}
\end{center}

\small $^*$En $T=1$, la coincidencia es casual (pocos datos).

\vspace{0.2cm}
\normalsize
Observe cómo $\bar{W}_D(T)$ es el que más tarda en converger: los primeros clientes tenían tiempos artificialmente bajos.
\end{frame}

% ============================== SLIDE 37 ==============================
\begin{frame}{Interpretación geométrica de la identidad de Little}

\begin{center}
\begin{tikzpicture}[scale=0.85]
    % Ejes
    \draw[->] (0,0) -- (12,0) node[right] {$t$};
    \draw[->] (0,0) -- (0,5) node[above] {Clientes};
    
    % Llegadas acumuladas A(t)
    \draw[thick, azuloscuro] (0,0) -- (1,1) -- (1.5,2) -- (2.5,3) -- (3,4) -- (4,5) -- (5.5,6) -- (7,7);
    \node[azuloscuro, right] at (7.1,7) {$A(t)$};
    
    % Partidas acumuladas D(t)
    \draw[thick, rojoclaro] (0,0) -- (1.8,0) -- (2.5,1) -- (3,2) -- (4,3) -- (5,4) -- (6,5) -- (7.5,6) -- (9,7);
    \node[rojoclaro, right] at (9.1,7) {$D(t)$};
    
    % Área entre curvas = integral de N(t) = suma de W_i
    \fill[blue!10, opacity=0.5] (0,0) -- (1,1) -- (1.5,2) -- (2.5,3) -- (3,4) -- (4,5) -- (5.5,6) -- (7,7) -- (9,7) -- (7.5,6) -- (6,5) -- (5,4) -- (4,3) -- (3,2) -- (2.5,1) -- (1.8,0) -- cycle;
    
    % Etiquetas
    \node at (4.5,3.5) {\textbf{Área} $= \displaystyle\int_0^T N(t)\,dt$};
    \node at (4.5,2.5) {$= \displaystyle\sum_i W_i^*$};
    
    % N(t) = A(t) - D(t)
    \draw[<->, thick, verdeoscuro] (4,5) -- (4,3);
    \node[verdeoscuro, right] at (4.1,4) {$N(t)$};
    
\end{tikzpicture}
\end{center}

\vspace{0.2cm}
El área entre $A(t)$ y $D(t)$ se puede leer horizontalmente ($\sum W_i$) o verticalmente ($\int N(t)\,dt$). Esta dualidad \textbf{es} la Ley de Little.
\end{frame}

% ============================== SLIDE 38 ==============================
\begin{frame}{Resumen: cuándo confiar en las métricas simuladas}

\begin{center}
\begin{tikzpicture}[
    box/.style={draw, rounded corners, text width=4.8cm, minimum height=2.2cm, align=left, font=\small},
]
    \node[box, fill=red!10] (trans) at (0,0) {
        \textbf{\textcolor{rojoclaro}{En el transitorio:}}\\[0.15cm]
        $\bullet$ $\bar{L}(T) \neq \lambda \cdot \bar{W}_D(T)$\\
        $\bullet$ Las métricas están sesgadas\\
        $\bullet$ Los IC no son válidos\\
        $\bullet$ Solución: truncar warm-up
    };
    
    \node[box, fill=green!10] (estable) at (7,0) {
        \textbf{\textcolor{verdeoscuro}{En estado estable:}}\\[0.15cm]
        $\bullet$ $L = \lambda \cdot W$ se cumple\\
        $\bullet$ Las métricas convergen\\
        $\bullet$ Los IC son confiables\\
        $\bullet$ Todas las versiones coinciden
    };
    
    \node[box, fill=blue!10] (siempre) at (3.5,-3.5) {
        \textbf{\textcolor{azuloscuro}{Siempre cierto:}}\\[0.15cm]
        $\bullet$ $\bar{L}(T) = \bar{\lambda}(T) \cdot \bar{W}_A(T)$\\
        $\bullet$ (Identidad contable)\\
        $\bullet$ Incluye tiempos parciales\\
        $\bullet$ No requiere estado estable
    };
    
    \draw[->, thick] (trans) -- (estable) node[midway, above] {$T_0$};
    \draw[->, thick, dashed] (siempre) -- (trans);
    \draw[->, thick, dashed] (siempre) -- (estable);
\end{tikzpicture}
\end{center}
\end{frame}

% ======================================================================
\section{Extensiones y Buenas Prácticas}
% ======================================================================

% ============================== SLIDE 39 ==============================
\begin{frame}{Sensibilidad a $\rho$: impacto en el transitorio}

\textbf{A mayor $\rho$, más largo es el transitorio:}

\begin{center}
\begin{tikzpicture}
\begin{axis}[
    width=13cm, height=5.5cm,
    xlabel={Tiempo de simulación},
    ylabel={$\bar{L}(T) / L^*$},
    xmin=0, xmax=150,
    ymin=0, ymax=1.3,
    grid=both, grid style={gray!20},
    legend pos=south east,
]
% rho = 0.5
\addplot[thick, verdeoscuro, domain=0:150, samples=200] 
    {1 - exp(-0.15*x)};
\addlegendentry{$\rho = 0.50$}

% rho = 0.83
\addplot[thick, azuloscuro, domain=0:150, samples=200] 
    {1 - exp(-0.06*x)};
\addlegendentry{$\rho = 0.83$}

% rho = 0.95
\addplot[thick, rojoclaro, domain=0:150, samples=200] 
    {1 - exp(-0.02*x)};
\addlegendentry{$\rho = 0.95$}

\addplot[dashed, black, domain=0:150] {1};

% T0 markers
\draw[dashed, verdeoscuro] (axis cs:20,0) -- (axis cs:20,1.3);
\draw[dashed, azuloscuro] (axis cs:50,0) -- (axis cs:50,1.3);
\draw[dashed, rojoclaro] (axis cs:120,0) -- (axis cs:120,1.3);

\end{axis}
\end{tikzpicture}
\end{center}

\vspace{0.1cm}
Para $\rho = 0.95$, el transitorio puede durar \textbf{5--10 veces más} que para $\rho = 0.50$. Sistemas muy cargados requieren simulaciones mucho más largas.
\end{frame}

% ============================== SLIDE 40 ==============================
\begin{frame}{Buenas prácticas en simulación de colas}

\textbf{Checklist para un estudio de simulación robusto:}

\vspace{0.3cm}
\begin{enumerate}
    \item \textbf{Verificar estabilidad}: Asegurar $\rho < 1$ antes de buscar estado estable.
    
    \vspace{0.15cm}
    \item \textbf{Detectar el transitorio}: Usar Welch u otro método. No asumir un $T_0$ arbitrario.
    
    \vspace{0.15cm}
    \item \textbf{Verificar Little}: Comprobar $\hat{L} \approx \hat{\lambda}\cdot\hat{W}$ como diagnóstico de que el sistema alcanzó estabilidad.
    
    \vspace{0.15cm}
    \item \textbf{Usar suficientes réplicas}: Mínimo $R \geq 10$ para intervalos de confianza.
    
    \vspace{0.15cm}
    \item \textbf{Corridas suficientemente largas}: $T \gg T_0$ (al menos $T/T_0 > 10$).
    
    \vspace{0.15cm}
    \item \textbf{Validar contra lo conocido}: Si existe solución analítica, comparar.
    
    \vspace{0.15cm}
    \item \textbf{Reportar incertidumbre}: Siempre acompañar estimaciones con intervalos de confianza.
\end{enumerate}
\end{frame}

% ============================== SLIDE 41 ==============================
\begin{frame}{¿Qué pasa si Little no se cumple en mi simulación?}

\textbf{Diagnóstico:} si $\hat{L} \neq \hat{\lambda}\cdot\hat{W}$ significativamente...

\vspace{0.4cm}
\begin{columns}
\column{0.48\textwidth}
\textcolor{rojoclaro}{\textbf{Posibles causas:}}
\begin{itemize}
    \item Período transitorio no eliminado
    \item Simulación demasiado corta
    \item Bug en el código (contadores mal actualizados)
    \item $\bar{W}$ calculado solo con completados
    \item Sistema inestable ($\rho \geq 1$)
\end{itemize}

\column{0.48\textwidth}
\textcolor{verdeoscuro}{\textbf{Acciones correctivas:}}
\begin{itemize}
    \item Aumentar $T_0$ (warm-up)
    \item Aumentar $T$ (largo de corrida)
    \item Revisar lógica de eventos
    \item Usar métrica consistente
    \item Verificar $\hat{\rho}$ observado
\end{itemize}
\end{columns}

\vspace{0.4cm}
\begin{block}{Little como herramienta de verificación}
La Ley de Little es una de las herramientas más poderosas para \textbf{validar} un simulador. Si no se cumple, algo está mal en la simulación o en la recolección de datos.
\end{block}
\end{frame}

% ======================================================================
\section{Conclusiones}
% ======================================================================

% ============================== SLIDE 42 ==============================
\begin{frame}{Conclusiones principales}

\begin{enumerate}
    \item La \textbf{simulación de eventos discretos} permite estimar métricas de colas cuando no hay soluciones analíticas.
    
    \vspace{0.2cm}
    \item La \textbf{Ley de Little} $L = \lambda W$ es fundamental pero requiere estado estable en su forma estocástica.
    
    \vspace{0.2cm}
    \item La \textbf{versión operacional completa} $\bar{L}(T) = \bar{\lambda}(T)\cdot\bar{W}_A(T)$ es una identidad que siempre se cumple, incluso en el transitorio.
    
    \vspace{0.2cm}
    \item La versión con $\bar{W}_D$ (solo completados) \textbf{no coincide} con $\bar{L}$ durante el transitorio debido al sesgo de selección.
    
    \vspace{0.2cm}
    \item El \textbf{período transitorio} debe detectarse y eliminarse para obtener estimaciones insesgadas del estado estable.
    
    \vspace{0.2cm}
    \item Para nuestro $M/M/2$ con $\rho = 0.833$, la simulación con warm-up de 50 minutos reproduce las métricas teóricas con alta precisión.
\end{enumerate}
\end{frame}

% ============================== SLIDE 43 ==============================
\begin{frame}{Resumen visual}

\begin{center}
\begin{tikzpicture}[>=Stealth,
    cbox/.style={draw, rounded corners, text width=3cm, minimum height=1.5cm, align=center, font=\small\bfseries},
]
    \node[cbox, fill=blue!15] (sim) at (0,3) {Simulación\\Eventos Discretos};
    \node[cbox, fill=orange!15] (trans) at (-5,0) {Período\\Transitorio};
    \node[cbox, fill=green!15] (ee) at (5,0) {Estado\\Estable};
    \node[cbox, fill=red!15] (little) at (0,0) {Ley de\\Little};
    \node[cbox, fill=purple!15] (met) at (0,-3) {Métricas:\\$L, W, L_q, W_q, \rho$};
    
    \draw[->, thick] (sim) -- (trans) node[midway, above left, font=\scriptsize] {warm-up};
    \draw[->, thick] (sim) -- (ee) node[midway, above right, font=\scriptsize] {largo plazo};
    \draw[->, thick] (trans) -- (little) node[midway, left, font=\scriptsize] {identidad contable};
    \draw[->, thick] (ee) -- (little) node[midway, right, font=\scriptsize] {$L=\lambda W$};
    \draw[->, thick] (little) -- (met) node[midway, right, font=\scriptsize] {verificaci\'{o}n};
    \draw[->, thick] (sim) -- (met) node[midway, right, font=\scriptsize] {estimaci\'{o}n};
    
\end{tikzpicture}
\end{center}
\end{frame}

% ============================== SLIDE 44 ==============================
\begin{frame}{Referencias y lecturas recomendadas}

\begin{thebibliography}{9}
\small

\bibitem{law} Law, A.M. (2015). \textit{Simulation Modeling and Analysis}, 5th ed. McGraw-Hill.

\bibitem{banks} Banks, J., Carson, J.S., Nelson, B.L., Nicol, D.M. (2014). \textit{Discrete-Event System Simulation}, 5th ed. Pearson.

\bibitem{little} Little, J.D.C. (1961). A proof for the queuing formula: $L = \lambda W$. \textit{Operations Research}, 9(3), 383--387.

\bibitem{stidham} Stidham, S. (1974). A last word on $L = \lambda W$. \textit{Operations Research}, 22(2), 417--421.

\bibitem{whitt} Whitt, W. (1991). A review of $L = \lambda W$ and extensions. \textit{Queueing Systems}, 9, 235--268.

\bibitem{welch} Welch, P.D. (1983). The statistical analysis of simulation results. In \textit{The Computer Performance Modeling Handbook}, Academic Press.

\bibitem{ross} Ross, S.M. (2014). \textit{Introduction to Probability Models}, 11th ed. Academic Press.

\end{thebibliography}
\end{frame}

% ============================== SLIDE 45 ==============================
\begin{frame}{}

\begin{center}
\vspace{1cm}
{\Huge\textbf{¿Preguntas?}}

\vspace{1cm}
{\large Simulación de Eventos Discretos}

\vspace{0.5cm}
{\large\textit{``Little's Law is the most useful and general result\\in queueing theory.''} --- S. Stidham}


\end{center}
\end{frame}

\end{document}
